\documentclass[11pt]{article}
\usepackage[a4paper,margin=25mm]{geometry}
\usepackage{amsmath,amssymb,amsthm,mathtools}
\usepackage{microtype,booktabs,hyperref}
\hypersetup{colorlinks=true,linkcolor=blue,citecolor=blue,urlcolor=blue}
\newcommand{\norm}[1]{\left\lVert#1\right\rVert}
\newcommand{\ip}[2]{\langle#1,#2\rangle}
\newcommand{\supp}{\operatorname{supp}}
\newcommand{\dist}{\operatorname{dist}}
\newcommand{\zer}{\operatorname{zer}}
\newcommand{\prox}{\operatorname{prox}}
\newcommand{\Id}{\operatorname{Id}}
\newtheorem{theorem}{Theorem}[section]
\newtheorem{proposition}[theorem]{Proposition}
\newtheorem{lemma}[theorem]{Lemma}
\newtheorem{corollary}[theorem]{Corollary}
\theoremstyle{remark}\newtheorem{remark}[theorem]{Remark}
\title{A sharp automatic Hessian threshold and complete proximal dynamics\\for least-squares norm-difference regularization}
\author{PPA research project}
\date{October 10, 2026}
\begin{document}
\maketitle
\begin{abstract}
For least-squares regularization by $\norm{x}_1-\norm{x}_p$, we prove
that every nonzero local minimum has a positive definite active Hessian
when $1<p\le3/2$, assuming only that the sensing matrix has no zero
columns. A moment identity excludes null directions below the threshold;
at its endpoint the quartic equality case is excluded by a negative
sixth-order objective term. For every finite $p>3/2$, an invertible
 two-dimensional example has a strict local minimum with singular active
Hessian. Strict complementarity and support independence have wider
ranges and are stated separately. An explicit common radius and positive
step interval certify the complete limiting-subdifferential resolvent,
including exclusion of every remote output, equality with the global
proximal map, identification, invariance, and variable-step local linear
convergence. For a specified $p=7/4$ example we obtain the sharp true
residual exponent $1/3$ and an exact $k^{-1/2}$ asymptotic for the actual
complete proximal iteration, with a separate geometric law on the radial
axis. No invariance of the radial critical curve is assumed.
\end{abstract}

\section{Model, scope, and relation to earlier work}
Let $A\in\mathbb R^{m\times n}$, $b\in\mathbb R^m$, and $\eta>0$.
Throughout the structural statements, every column of $A$ is nonzero.
For $p>1$ set
\begin{equation}\label{eq:model}
 \Phi_p(x)=\tfrac12\norm{Ax-b}_2^2+
 \eta(\norm{x}_1-\norm{x}_p),\qquad F_p=\partial_{\rm lim}\Phi_p,
 \qquad Q=A^TA,\quad c=A^Tb.
\end{equation}
Here $\partial_{\rm lim}$ denotes the limiting subdifferential. The
least-squares compatibility $c=A^Tb$ is retained, including for singular
$Q$. For every $\lambda>0$ we use the complete relation
\begin{equation}\label{eq:complete}
 J_{\lambda F_p}(z)=\{y\in\mathbb R^n:z-y\in\lambda F_p(y)\},
 \qquad r_{F_p}(x)=\dist(0,F_p(x)).
\end{equation}
Its outputs are unrestricted until an exclusion argument localizes them.
The global proximal map is the set of global minimizers of
$\Phi_p(y)+\norm{y-z}^2/(2\lambda)$; for a nonconvex function it need
not equal \eqref{eq:complete} without proof. All norms without subscripts
on points, matrices, or gradients are Euclidean norms or their operator
norms.

The norm-difference model is established prior work. Wang's penalty
formulation includes $\ell_1-\epsilon\ell_r$ for $r>1$\cite{WangPenalty}.
Wang--Zhang study $1<p\le2$ and a smoothed neural optimization
method\cite{WangZhang}; Huo--Chen--Ge--Ng study
$\ell_1-\beta\ell_q$ and recovery estimates\cite{Huo}.
Support-column independence for the $p=2$ constrained and penalized models
is already given by Yin--Lou--He--Xin, Theorems~2.3--2.4\cite{Yin};
we do not present support independence as a new standalone principle.
Lou--Yan derive a proximal operator and analyze related algorithms for
$p=2$\cite{LouYan}. The proposed model-specific contribution is the
universal automatic active-Hessian threshold and its endpoint mechanism.
The quantitative complete-resolvent conclusions are consequences of
that structure; ordinary local strong-monotonicity and analytic/KL
methods remain applicable where their conditions hold.

The fractional framework of Zhao--He--Ma--Wang\cite{Fractional}
relates constrained ratio and norm-difference stationarity (Theorem~3.1)
and supplies conditional KL rates (Theorem~6.4). This first-order
equivalence includes a point-dependent norm coefficient equal to one;
it does not itself prove automatic active curvature for the penalized
least-squares model. A conditional KL rate likewise does not supply
the universal curvature threshold.

\paragraph{Priority boundary.}
The original published 2015 support-independence proofs have been read.
For Wang--Zhang (2017), the publisher abstract, reference list, and
Appendix proofs were available, but the main body and theorem statements
were not. For Huo et al.\ (2023), only the publisher abstract/reference
list and institutional metadata were available. The relevant
Zhao--He--Huang--Huang norm-difference neural-network paper\cite{ZhaoNeural}
was available only as a preview. Accordingly this paper
makes no claim of certified first publication of the threshold. The
mathematical statements below do not depend on any inaccessible body.
The separate source audit records the precise access boundary and
comparison obligations; inability to retrieve a source is not evidence
that it contains no matching theorem.

\section{Automatic structure and the sharp threshold}
\begin{theorem}\label{thm:automatic}
For \eqref{eq:model} with no zero columns of $A$, the following statements
hold at every nonzero local minimum $\bar x$, with $I=\supp\bar x$.
\begin{enumerate}
\item For $1<p<2$, strict complementarity is automatic:
$|(Q\bar x-c)_j|<\eta$ for every $j\notin I$.
\item For every $p>1$, $A_I$ has linearly independent columns.
\item For $1<p\le3/2$, the active Hessian
\begin{equation}\label{eq:activeH}
 H_I=Q_{II}-\eta\nabla^2\norm{\bar x_I}_p
\end{equation}
is positive definite. In particular $|I|\le\operatorname{rank}A$.
\end{enumerate}
The exponent $3/2$ in item~3 is sharp across all finite $p>1$: for each
finite $p>3/2$, there is a two-dimensional example with invertible $A$
and a nonzero strict local minimum whose active Hessian is singular.
\end{theorem}
\begin{proof}[Strict complementarity]
At a nonzero point the $p$-norm is $C^1$, including at zero coordinates,
and the exact smooth sum rule gives
\begin{equation}\label{eq:Fnonzero}
 F_p(x)=Qx-c+\eta(\partial\norm{x}_1-\nabla\norm{x}_p),\quad x\ne0.
\end{equation}
Fermat's rule yields $|(Q\bar x-c)_j|\le\eta$ off $I$.
If equality holds, choose $\sigma\in\{-1,1\}$ so that
$\sigma(Q\bar x-c)_j=-\eta$, and let $R=\norm{\bar x}_p>0$.
For $t>0$ the exact difference is
\[
 \Phi_p(\bar x+\sigma t e_j)-\Phi_p(\bar x)
 =\tfrac12\norm{A_j}^2t^2-
 \eta\big((R^p+t^p)^{1/p}-R\big).
\]
The second term is $-\eta R^{1-p}t^p/p+O(t^{2p})$.
For $p<2$ this is negative for all sufficiently small positive $t$,
contradicting local minimality.
\end{proof}
\begin{proof}[Support independence]
On the fixed active face the objective is analytic and the local-minimum
second-order condition gives $H_I\succeq0$. At a vector all of whose
active coordinates are nonzero, the Hessian of the $p$-norm is positive
semidefinite, with kernel precisely the radial line. Indeed, for
$R=\norm{u}_p$, $w_i=|u_i|^p/R^p$, and $r_i=z_i/u_i$,
\begin{equation}\label{eq:normkernel}
 z^T\nabla^2\norm{u}_p z
 =(p-1)R\left(\sum_iw_i r_i^2-
                    \left(\sum_iw_i r_i\right)^2\right).
\end{equation}
If $A_Iz=0$ for nonzero $z$, positivity of $H_I$ and
\eqref{eq:normkernel} force $z=a\bar x_I$ for $a\ne0$; hence
$A\bar x=0$. Dot active stationarity with $\bar x$:
\begin{equation}\label{eq:radialcompat}
 0=\norm{A\bar x}^2-b^TA\bar x+
                 \eta(\norm{\bar x}_1-\norm{\bar x}_p).
\end{equation}
For $|I|\ge2$ the last term is positive. For $|I|=1$,
$A\bar x=0$ contradicts the nonzero column assumption. Both cases
exclude $z$. This argument holds for every $p>1$.
\end{proof}
\begin{proof}[Positive definiteness up to the endpoint]
Suppose that $H_I\succeq0$ is singular and choose nonzero
$z\in\ker H_I$. Apply an orthogonal sign change, assume the active coordinates are positive, and put
$x_i=|\bar x_i|>0$, $R=(\sum_i x_i^p)^{1/p}$,
\[
 w_i=x_i^p/R^p,\quad r_i=z_i/x_i,\quad
 a=\sum_iw_i r_i,\quad s_i=r_i-a,\quad
 \mu_k=\sum_iw_i s_i^k.
\]
Thus $\mu_1=0$. If $\mu_2=0$, $z$ is radial; norm homogeneity and
$H_Iz=0$ imply $A\bar x=0$, contradicting
\eqref{eq:radialcompat} as above. Hence $\mu_2>0$.
Factor the radial part exactly, for sufficiently small real $t$:
\[
 \frac{\norm{\bar x+tz}_p}{R}
 =(1+at)f\left(\frac{t}{1+at}\right),\qquad
 f(u)=\left(\sum_iw_i(1+us_i)^p\right)^{1/p}.
\]
Writing $f(u)=1+c_2u^2+c_3u^3+c_4u^4+O(u^5)$ gives
\begin{align}\label{eq:coeffs}
 c_2&=\tfrac12(p-1)\mu_2,&
 c_3&=\tfrac16(p-1)(p-2)\mu_3,\nonumber\\
 c_4&=\tfrac1{24}(p-1)(p-2)(p-3)\mu_4
                         -\tfrac18(p-1)^3\mu_2^2.
\end{align}
Along the original line the cubic coefficient of the norm is
$c_3-ac_2$. The objective's linear term vanishes by stationarity and
its quadratic term vanishes since $z\in\ker H_I$. All coefficients
from degree three onwards are minus $\eta R$ times the corresponding
norm coefficients. Local minimality on both sides therefore requires
$a=c_3/c_2$, and the norm's quartic coefficient becomes
$C_4=c_4-c_3^2/c_2$.

The following exact square identity supplies its sign:
\begin{equation}\label{eq:pearson}
 \mu_4-\mu_2^2-\frac{\mu_3^2}{\mu_2}
 =\sum_iw_i\left(s_i^2-\mu_2-\frac{\mu_3}{\mu_2}s_i\right)^2\ge0.
\end{equation}
Substitution in \eqref{eq:coeffs} yields
\begin{equation}\label{eq:C4}
 C_4\ge(p-1)\left(
 \frac{(3-2p)(p+1)}{24}\mu_2^2+
 \frac{(2-p)(p+1)}{72}\frac{\mu_3^2}{\mu_2}\right).
\end{equation}
For $1<p<3/2$ it is strictly positive, so the objective has a negative
quartic term, a contradiction.

At $p=3/2$ the right side is $5\mu_3^2/(576\mu_2)$.
If $C_4>0$, the same contradiction applies. Otherwise
$\mu_3=0$, $\mu_4=\mu_2^2$, and \eqref{eq:pearson} forces
$s_i^2=\mu_2$ at every coordinate. The values $s_i=\pm d$,
$d=\sqrt{\mu_2}>0$, have total weight $1/2$ each, and $a=0$.
The exact norm along the line is consequently
\begin{align}\label{eq:endpoint}
 \frac{\norm{\bar x+tz}_p}{R}
 &=\left(\frac{(1+dt)^{3/2}+(1-dt)^{3/2}}2\right)^{2/3}\nonumber\\
 &=1+\frac{d^2t^2}4+\frac{d^6t^6}{192}+O(t^8).
\end{align}
Its cubic, quartic, and fifth-order terms vanish. After the quadratic
cancellation, the objective has strictly negative sixth-order coefficient
$-\eta R d^6/192$, again impossible at a local minimum. Thus $H_I\succ0$.
\end{proof}

\begin{remark}[A second proof by radial relaxation]
The endpoint sixth-order certificate above is retained. A different
fourth-order argument explains its geometric mechanism. For a hypothetical
active null direction $z$, write
$\norm{\bar x+tz}_p=R+n_1t+n_2t^2+n_4t^4+O(t^5)$;
the cubic vanishes by local minimality and $n_2>0$. Put
$q_0=\norm{A\bar x}^2>0$. Homogeneity and $H_Iz=0$ give
$\bar x^TQz=0$, so for $x=(1+s)\bar x+tz$,
\[
 \Phi_p(x)-\Phi_p(\bar x)=\frac{q_0s^2}{2}
 +\eta n_2\frac{st^2}{1+s}
 -\eta n_4\frac{t^4}{(1+s)^3}+O(t^5).
\]
Choosing $s=-\eta n_2t^2/q_0$ gives quartic coefficient
$-\eta n_4-\eta^2n_2^2/(2q_0)$. Thus a singular active minimum would
require $n_4\le-\eta n_2^2/(2q_0)<0$. The Pearson bound gives
$n_4=RC_4\ge0$ for $1<p\le3/2$, a contradiction. The perturbation
preserves active signs and leaves inactive coordinates zero.
\end{remark}

\begin{proposition}\label{prop:sharpfamily}
For each finite $p>3/2$, let $R=2^{1/p}$,
$\beta=R(p-1)/2$, $c_4=-(p-1)(2p-3)(p+1)/24<0$, and choose
\[
 \alpha>\alpha_0:=\frac{\beta^2}{-4Rc_4}
       =\frac{3R(p-1)}{2(2p-3)(p+1)}.
\]
Let $Q\succ0$ have eigenvalues $\alpha$ in direction $(1,1)$ and
$\beta$ in direction $(1,-1)$, and put
$A=Q^{1/2}$, $\bar x=(1,1)$,
$c=Q\bar x+(1-R/2)(1,1)$, $b=A^{-T}c$, and $\eta=1$.
Then $\bar x$ is a strict local minimum, with singular active Hessian.
It is an isolated stationary point. No local true-residual error bound
at $\bar x$ has exponent greater than $1/3$.
\end{proposition}
\begin{proof}
Stationarity is immediate; the norm Hessian has radial eigenvalue zero
and tangent eigenvalue $\beta$. Thus the active Hessian eigenvalues are
$\alpha$ and zero. Write $x=(1+h+t,1+h-t)$ and let
$V(h,t)=\Phi(x)-\Phi(\bar x)$. On the positive-coordinate
neighborhood the norm is analytic for every finite real $p>1$. Expansion gives
\begin{equation}\label{eq:generalquartic}
 V(h,t)=\alpha h^2+\beta ht^2-Rc_4t^4+
            O(h^2t^2+|h|t^4+t^6).
\end{equation}
The analytic implicit function theorem solves $V_h=0$ as an even
curve $h=g(t)=-\beta t^2/(2\alpha)+O(t^4)$. Its reduced objective is
\[
 V(g(t),t)=Kt^4+O(t^6),\qquad K=-Rc_4-\beta^2/(4\alpha)>0.
\]
For $\alpha=2\alpha_0$, $K=R(p-1)(2p-3)(p+1)/48>0$.
Shrink the neighborhood so that $V_{hh}\ge\alpha$ and the reduced
objective is at least $Kt^4/2$. Taylor's integral formula gives
$V(h,t)\ge\alpha|h-g(t)|^2/2+Kt^4/2$, proving full-space strictness.
Every nearby stationary point lies on $g$, where the reduced derivative
is $4Kt^3+O(t^5)$; hence only $t=0$ is stationary after shrinking.
Along $g$, $\norm{x-\bar x}\sim\sqrt2|t|$ while
$r_F(x)=|V_t(g(t),t)|/\sqrt2=O(|t|^3)$.
If $q>1/3$, the quotient $\norm{x-\bar x}/r_F(x)^q$ diverges.
Stationary isolation makes the same statement valid for distance to the
entire stationary set in a sufficiently small ball. This proof gives no
uniform exponent for all models with $p>3/2$.
\end{proof}

\begin{remark}[Zero columns, the origin, and other stationary points]
If $A=[0,1]$, $b=0$, and $\eta=1$, every $(a,0)$, $a\ne0$, is a
local minimum with zero active Hessian: on its active ray the objective
is constant, and off the ray the positive $|y|$ term dominates the
negative $O(|y|^p)$ correction. This is a material exception.
At the origin, examination of both directions on every coordinate axis
shows that local minimality requires $A^Tb=0$. Conversely, if $A^Tb=0$
and $A$ has no zero columns, the origin is the unique global minimum:
the norm-difference is nonnegative, and it vanishes only at one-sparse
vectors, where the data quadratic is positive. Dotting stationarity
with a nonzero $x$ excludes every other stationary point. Formula
\eqref{eq:Fnonzero} must not be used at zero. The theorem's quantifier is
local minima, not arbitrary stationary points or arbitrary PPA limits.
For example at $p=3/2$ put $R=2^{2/3}$, $d=R/8$, $\alpha>0$,
\[
 Q=\begin{pmatrix}d&-d\\-d&\alpha+d\end{pmatrix},\qquad
 c=(1-R/2,\alpha+1-R/2),\qquad \bar x=(1,1).
\]
Taking $A=Q^{1/2}$ and $b=A^{-T}c$ gives stationarity and
$H=\operatorname{diag}(0,\alpha)$, but
$\Phi(\bar x+te_1)-\Phi(\bar x)=Rt^3/32+O(t^4)$.
It is a stationary nonminimum.
\end{remark}

\begin{remark}[The separate strict-complementarity boundary]
For every finite $p\ge2$, take $\eta=1$, $Q=\operatorname{diag}(1,2)$,
$c=(1,1)$, $A=\operatorname{diag}(1,\sqrt2)$,
$b=(1,1/\sqrt2)$, and $\bar x=(1,0)$. If $r=\norm x_2$, then
\[
 \Phi_2(x)-\Phi_2(\bar x)=\tfrac12(r-1)^2+\tfrac12x_2^2
 +(|x_1|-x_1)+(|x_2|-x_2)\ge0,
\]
with equality only at $\bar x$. Since $\norm x_p\le\norm x_2$ for
$p\ge2$, the same point is the unique global minimum of $\Phi_p$.
Yet $|(Q\bar x-c)_2|=\eta$. Thus the automatic strict-complementarity
range is $1<p<2$, separately from the active-PD threshold $3/2$.
\end{remark}

\section{An explicit certificate for the complete resolvent}
Fix any nonzero local minimum in Theorem~\ref{thm:automatic} with
$1<p\le3/2$. Let $R=\norm{\bar x}_p$, $a_*=\min_{i\in I}|\bar x_i|$,
$s=|I|$, $C_n=n^{1/p-1/2}$, $C_s=s^{1/p-1/2}$, and
$h_* =\lambda_{\min}(H_I)>0$. If $I\ne\{1,\ldots,n\}$ set
\[
 \delta=\min_{j\notin I}\{\eta-|(Q\bar x-c)_j|\}>0,
 \qquad B_Q=\max_{j\notin I}\norm{Q_j}.
\]
Every condition containing $\delta$ or $B_Q$ is omitted for full support.
Put
\[
 r_0=\min\{a_*/2,R/(2C_n)\},\quad d=a_*/2,\quad N_0=R/2,
 \quad U=\max_{i\in I}|\bar x_i|+r_0,\quad W=\sqrt{s}\,U^{p-1}.
\]
For $s\ge2$, define the finite explicit bound
\begin{align}\label{eq:LH}
 L_H=(p-1)\big[&N_0^{1-p}(2-p)d^{p-3}
 +(p-1)N_0^{-p}C_s d^{p-2}\nonumber\\
 &+2N_0^{1-2p}W(p-1)d^{p-2}
 +(2p-1)N_0^{-2p}C_s W^2\big].
\end{align}
Choose any positive $\rho$ not exceeding each applicable member of
\begin{equation}\label{eq:rho}
 r_0,\qquad \frac{h_*}{2\eta L_H},\qquad
 \frac{\delta}{4B_Q},\qquad
 \frac R2\left(\frac{\delta}{4\eta}\right)^{1/(p-1)}.
\end{equation}
The $B_Q$ quotient is omitted if $B_Q=0$; the $L_H$ quotient is omitted
for $s=1$. Set $m=h_*/2$ for $s\ge2$, and
$m=\norm{A_I}^2$ for $s=1$. Finally set
\begin{equation}\label{eq:M}
 L_p=\sqrt n+C_n,\qquad
 M=\norm{Q\bar x-c}+\norm Q\rho/4+\eta L_p.
\end{equation}

\begin{theorem}\label{thm:fullprox}
For any $0<\lambda_-\le\lambda_+$ and $r>0$ satisfying
\begin{equation}\label{eq:windows}
 \lambda_+ M<\rho/2,\qquad
 r<\min\{\rho/4,\lambda_-\delta/2\},
\end{equation}
with the second entry omitted for full support, every
$\lambda\in[\lambda_-,\lambda_+]$ and every $z\in B_r(\bar x)$ have
exactly one complete resolvent output. It equals the global proximal
map, has support $I$ and the signs of $\bar x_I$, and obeys
\begin{equation}\label{eq:linear}
 \norm{J_{\lambda F_p}(z)-\bar x}
 \le(1+\lambda m)^{-1}\norm{z-\bar x}<r.
\end{equation}
For any $\lambda_k\in[\lambda_-,\lambda_+]$ and $x_0\in B_r(\bar x)$,
the complete PPA exists uniquely and
\[
 \norm{x_k-\bar x}\le(1+\lambda_-m)^{-k}\norm{x_0-\bar x}.
\]
In $B_\rho(\bar x)$ the strong true-residual error bound is
\begin{equation}\label{eq:linearEB}
 \norm{x-\bar x}\le\max\{1/m,2\rho/\delta\}\,r_{F_p}(x),
\end{equation}
with the second constant omitted for full support. In particular
$\bar x$ is an isolated stationary point and a strict local minimum.
All constants are computable for the specified data and minimum;
none is uniform over all models or minima.
\end{theorem}
\begin{proof}[Continuity and curvature budgets]
For $y\in B_\rho(\bar x)$, $\norm y_p\ge R-C_n\rho\ge R/2$.
Signs of active coordinates are preserved. Off support,
\[
 |(Qy-c)_j-\eta\nabla_j\norm y_p|
 \le |(Q\bar x-c)_j|+B_Q\rho+
                           \eta(2\rho/R)^{p-1}\le\eta-\delta/2.
\]
For active vectors $u$, with $v_i=\operatorname{sign}(u_i)|u_i|^{p-1}$,
\[
 \nabla^2\norm u_p=(p-1)
 \big[\norm u_p^{1-p}\operatorname{diag}(|u_i|^{p-2})
                  -\norm u_p^{1-2p}vv^T\big].
\]
On the active $r_0$ ball, the diagonal-power difference is bounded by
$(2-p)d^{p-3}\norm{u-\tilde u}$, the vector difference by
$(p-1)d^{p-2}\norm{u-\tilde u}$, and $\norm v\le W$.
The norm powers have mean-value constants
$(p-1)N_0^{-p}C_s$ and $(2p-1)N_0^{-2p}C_s$.
Apply these bounds to the two products, with
$\norm{vv^T-\tilde v\tilde v^T}\le2W\norm{v-\tilde v}$.
This gives exactly \eqref{eq:LH}, hence
$H_I(u)\succeq m\Id$ by \eqref{eq:rho}. For $s=1$ the active norm is
linear and its Hessian is identically zero.
\end{proof}
\begin{proof}[All outputs, independent existence, and dynamics]
The penalty $g_p=\norm{\cdot}_1-\norm{\cdot}_p$ is globally
$L_p$-Lipschitz. Every Fr\'echet subgradient has norm at most $L_p$:
combine its lower first-order bound in any unit direction with the
upper Lipschitz difference bound. Limiting subgradients inherit this
bound. Thus, including at $y=0$, the smooth sum rule gives
$F_p(y)=Qy-c+\eta\partial_{\rm lim}g_p(y)$ and any complete output
satisfies $z=y+\lambda(Qy-c+\eta w)$ with $\norm w\le L_p$.
Consequently
\begin{align}\label{eq:remote}
 (\Id+\lambda Q)(y-z)&=\lambda(c-Qz-\eta w),\nonumber\\
 \norm{y-z}&\le\lambda(\norm{Qz-c}+\eta L_p)\le\lambda M<\rho/2,
\end{align}
because $Q\succeq0$ and $\norm{(\Id+\lambda Q)^{-1}}\le1$.
Every complete output is therefore in $B_{3\rho/4}(\bar x)$,
where \eqref{eq:Fnonzero} is valid. If $j\notin I$ and $y_j\ne0$,
the offsupport bound and the resolvent equation give
\[
 \operatorname{sign}(y_j)z_j\ge |y_j|+\lambda\delta/2
 >\lambda\delta/2\ge\lambda_-\delta/2>r,
\]
contradicting $|z_j|\le\norm{z-\bar x}<r$.
All outputs are identified on the active face. Strong monotonicity on
that face, compared to $\bar x$, gives \eqref{eq:linear}; comparison
of two outputs at the same input gives uniqueness. No existence was
presumed in either argument.

Independently, the global proximal objective is continuous and coercive:
both data loss and penalty are nonnegative and its added quadratic is
coercive. It attains a global minimum. Fermat's rule places every such
minimum in the complete resolvent. Thus the resolvent is nonempty and
the two maps equal the singleton already proved unique. Invariance and
the variable-step estimate follow by induction.

For the residual bound, a nonzero inactive coordinate forces every
subgradient's corresponding component to have modulus at least
$\delta/2$, giving \eqref{eq:linearEB}. On the active face all inactive
intervals contain zero; the true residual equals the norm of the active
gradient, which is at least $m\norm{x-\bar x}$ by strong monotonicity.
This also proves isolation. To prove strictness, take a local-minimum
neighborhood. An interior point with the same objective value would
itself be a local minimum and hence stationary, contradicting isolation.
\end{proof}

\begin{remark}[One completely specified positive interval]
For the computed $\rho,m,\delta,M$, one may choose
$\lambda_+=\rho/(4M)$, $\lambda_-=\lambda_+/2$, and
$r=\min\{\rho/8,\lambda_-\delta/4\}$, omitting the gap entry for full
support. These choices satisfy all strict inequalities in
\eqref{eq:windows}. No SOSC, residual error bound, or resolvent
single-valuedness is assumed in advance.
\end{remark}


\clearpage
\section{A quartic minimum: sharp true residual and complete PPA}
\begin{theorem}\label{thm:slow}

Fix

\[
p=7/4,\quad \eta=1,\quad R=2^{4/7},\quad
Q=\frac R{16}\begin{pmatrix}19&13\\13&19\end{pmatrix},\quad
\bar x=(1,1),\quad c=(1+3R/2)(1,1),
\]

and take \(A=Q^{1/2}\), \(b=A^{-T}c\). In particular \(A^TA=Q\), \(A^Tb=c\), and both columns of \(A\) are nonzero. Define on all of \(\mathbb R^2\)

\[
\Phi(x)=\tfrac12\|Ax-b\|_2^2+\|x\|_1-\|x\|_{7/4},\quad
F=\partial_{\rm lim}\Phi,\quad S=F^{-1}(0),\quad
r_F(x)=\inf_{v\in F(x)}\|v\|_2.
\]

The complete resolvent is \(J_{\lambda F}(z)=\{y:z-y\in\lambda F(y)\}\), with every output \(y\in\mathbb R^2\) initially allowed. Put

\[
\mathcal B=\{x=(1+h+t,1+h-t): |h|,|t|\le 1/32\},\qquad
\mathcal U=\overline B_{1/64}(\bar x).
\]

The following assertions hold.

\begin{enumerate}
\item \(\bar x\) is a strict local minimum and is the sole stationary point in \(\mathcal B\). The active Hessian at \(\bar x\) has eigenvalues \(2R\) and \(0\).
\item On \(\mathcal B\) the \textbf{strong true-residual error bound} is
   \[
   \|x-\bar x\|_2\le (512/R)^{1/3}r_F(x)^{1/3}.
   \]
   The exponent \(1/3\) is optimal: no exponent greater than \(1/3\), with any finite constant and any neighborhood of \(\bar x\), satisfies this bound. On \(\mathcal U\), \(d(x,S)=\|x-\bar x\|_2\).
\item For \textbf{every} \(0<\lambda\le 1/128\) and \textbf{every} \(z\in\mathcal U\), the complete \(J_{\lambda F}(z)\) is a singleton, equals the global minimizing proximal map
   \[
   \mathop{\rm argmin}_{y\in\mathbb R^2}
   \left\{\Phi(y)+\frac{\|y-z\|_2^2}{2\lambda}\right\},
   \]
   and belongs to \(\mathcal U\).
\item Fix any one \(\lambda\in(0,1/128]\), and run the ordinary complete PPA from \textbf{any} \(x_0\in\mathcal U\): \(x_{k+1}\in J_{\lambda F}(x_k)\). This trajectory is unique and converges to \(\bar x\). Writing \(x_k=(1+h_k+t_k,1+h_k-t_k)\), if \(t_0\ne0\), then \(t_k\) keeps its initial sign, and
   \[
   \frac{h_k}{t_k^2}\longrightarrow-\frac3{32},\qquad
   \sqrt{k}|t_k|\longrightarrow\sqrt{\frac{128}{13\lambda R}},\qquad
   \boxed{\sqrt{k}\|x_k-\bar x\|_2\longrightarrow
   \frac{16}{\sqrt{13\lambda R}}.}
   \]
   If \(t_0=0\), the exact separate formula is \(t_k=0\), \(h_k=(1+2\lambda R)^{-k}h_0\).

\end{enumerate}

The constants \(1/64\) and \(1/128\) below are certified sufficient values, not maximal basin or step bounds. The asymptotic assertions fix \(\lambda\); they do not assert this particular constant for an arbitrary variable-step sequence.

\end{theorem}
\subsection{Exact scalar formula and analytic remainder certificate}

Throughout \(\mathcal B\), the two coordinates of \(x\) are positive. Put \(s=1+h\), \(u=t/s\), and

\[
f(u)=\left(\frac{(1+u)^{7/4}+(1-u)^{7/4}}2\right)^{4/7},\qquad
\varphi(h,t)=\Phi(1+h+t,1+h-t)-\Phi(\bar x).
\]

The eigenvalues of \(Q\) are \(2R\) in the radial direction \((1,1)\) and \(3R/8\) in the tangent direction \((1,-1)\). Substitution, including \(c=A^Tb\), gives the exact identity

\[
\frac{\varphi(h,t)}R=2h^2+\frac38t^2-s(f(u)-1). \tag{D.1}\label{eq:dyn1}
\]

In particular \(\nabla\Phi(\bar x)=0\). The symmetric binomial expansion gives

\[
f(u)=1+\frac38u^2-\frac{11}{256}u^4+E(u). \tag{D.2}\label{eq:dyn2}
\]

Here, for \textbf{all real} \(|u|\le1/16\),

\[
|E(u)|\le5|u|^6,\quad |E'(u)|\le28|u|^5,\quad
|E''(u)|\le140|u|^4. \tag{D.3}\label{eq:dyn3}
\]

For completeness these bounds have an analytic certificate, rather than a sampled verification. Let \(G(u)=((1+u)^{7/4}+(1-u)^{7/4})/2\). Its even binomial coefficients from degree two onward are positive; after degree four their absolute values decrease. For complex \(|u|\le r=7/8\),

\[
|G(u)-1|\le \frac{21}{32}r^2+
\frac{35}{2048}\frac{r^4}{1-r^2}
=\frac{214375}{393216}<1.
\]

Therefore the analytic branch \(G^{4/7}\) with value one at zero is defined on a neighborhood of this closed disk and has modulus at most \(2\). If \(f=\sum a_{2j}u^{2j}\), the Cauchy coefficient estimate is \(|a_{2j}|\le2r^{-2j}\). After removal of the displayed terms through degree four, put \(v=(u/r)^2\le1/196\). Summing the majorant series and its two derivatives gives

\[
\begin{aligned}
|E|/|u|^6&\le 2r^{-6}/(1-v)<5,\\
|E'|/|u|^5&\le 2r^{-6}\left[6/(1-v)+2v/(1-v)^2\right]<28,\\
|E''|/|u|^4&\le 2r^{-6}\left[30/(1-v)+26v/(1-v)^2
                         +8v^2/(1-v)^3\right]<140.
\end{aligned}
\]

The inequalities at zero are interpreted by continuity. This proves \eqref{eq:dyn3} on the full interval. Direct expansion of \eqref{eq:dyn1} consequently gives

\[
\frac{\varphi}R=2h^2+\frac38ht^2+\frac{11}{256}t^4
+O(h^2t^2+|h|t^4+t^6). \tag{D.4}\label{eq:dyn4}
\]

\subsection{Star inequality, strictness, and the sharp residual exponent}

The quantity needed for invariance is the gradient inner product with the displacement, not the objective in \eqref{eq:dyn4}. Exactly,

\[
D(h,t):=\langle x-\bar x,\nabla\Phi(x)\rangle
=h\varphi_h+t\varphi_t.
\]

Differentiating \eqref{eq:dyn1} and using \eqref{eq:dyn2} gives

\[
\frac DR=4h^2+\frac98hu^2+\frac34h^2u^2
+\left(\frac{11}{64}+\frac{11h}{256}\right)u^4-hE(u)-uE'(u).
\tag{D.5}\label{eq:dyn5}
\]

On \(\mathcal B\), \(|u|\le1/31<1/16\), so \eqref{eq:dyn3} applies. Complete the first square in \eqref{eq:dyn5}, discard the nonnegative term \(3h^2u^2/4\), and bound the remaining terms:

\[
\frac DR\ge4\left(h+\frac9{64}u^2\right)^2+
\left[\frac{95}{1024}-\frac{11}{8192}
-\frac{28+5/32}{961}\right]u^4
\ge4\left(h+\frac9{64}u^2\right)^2+\frac1{32}u^4.
\tag{D.6}\label{eq:dyn6}
\]

The last rational bracket exceeds \(1/32\). Also \((33/32)^4<2\), so

\[
h^2+t^4\le2\left(h+\frac9{64}u^2\right)^2+
\left(2\left(\frac9{64}\right)^2+2\right)u^4.
\]

Comparison with \eqref{eq:dyn6} yields the conservative full-box bound

\[
\boxed{D(h,t)\ge\frac R{128}(h^2+t^4).} \tag{D.7}\label{eq:dyn7}
\]

One can sharpen \eqref{eq:dyn7} to \(R(h^2+t^4)/64\) by keeping \((33/32)^4\) rather than replacing it by two: the required rational check is
\((2(9/64)^2+(33/32)^4)/64<1/32\). We use this certified sharper bound below:

\[
D(h,t)\ge\frac R{64}(h^2+t^4). \tag{D.8}\label{eq:dyn8}
\]

Every ray segment from \(\bar x\) to a point of \(\mathcal B\) stays in \(\mathcal B\). Integrating \eqref{eq:dyn8} along that segment gives

\[
\varphi(h,t)\ge \frac R{128}h^2+\frac R{256}t^4>0
\quad\text{when }(h,t)\ne(0,0).
\]

This proves strict local minimality. A stationary point in the box has \(D=0\), so \eqref{eq:dyn8} proves stationarity isolation there. The Hessian eigenvalues stated in the theorem follow from \eqref{eq:dyn1} or from cancellation of the tangent norm Hessian against \(3R/8\).

Since \(\mathcal B\) has positive coordinates, \(F(x)=\{\nabla\Phi(x)\}\) there: this is the actual complete residual. Put \(d=\|x-\bar x\|_2=\sqrt{2(h^2+t^2)}\). Since \(|h|<1\),

\[
h^2+t^4\ge h^4+t^4\ge\tfrac12(h^2+t^2)^2=d^4/8.
\]

Consequently \(d\,r_F(x)\ge D(h,t)\ge Rd^4/512\), proving the strong exponent-\(1/3\) error bound, including \(d=0\). Every other stationary point is outside \(\mathcal B\), hence at distance greater than \(\sqrt2/32\) from \(\bar x\). For \(x\in\mathcal U\), its distance to such a point exceeds \(\sqrt2/32-1/64>1/64\ge d\), proving the asserted equality with distance to the entire \(S\).

To prove optimality, \(\varphi_{hh}(0,0)=4R>0\), and the analytic implicit function theorem gives the unique even radial critical curve near zero:

\[
\varphi_h(h_*(t),t)=0,\qquad
h_*(t)=-\frac3{32}t^2+O(t^4).
\]

Substitution into \eqref{eq:dyn4} gives

\[
\varphi(h_*(t),t)=\frac{13R}{512}t^4+O(t^6),\qquad
\varphi_t(h_*(t),t)=\frac{13R}{128}t^3+O(t^5).
\]

The exact gradient-coordinate transformation is
\(\|\nabla\Phi\|_2=\sqrt{\varphi_h^2+\varphi_t^2}/\sqrt2\).
Thus on this curve \(d\sim\sqrt2|t|\) whereas
\(r_F\sim13R|t|^3/(128\sqrt2)\). For any \(q>1/3\), \(d/r_F^q\) tends to infinity. The critical curve is used only to test sharpness; it is not declared invariant under PPA.

\subsection{Complete-fiber exclusion, global proximal equality, and invariance}

Let \(g(x)=\|x\|_1-\|x\|_{7/4}\). It is globally Euclidean-Lipschitz with constant
\(L_p=\sqrt2+2^{1/14}\). Every Fr\'echet subgradient has norm at most \(L_p\), by testing the subgradient inequality in unit directions against the Lipschitz upper bound; every limiting subgradient inherits this bound. The exact smooth-sum rule, valid including at zero, is

\[
F(y)=Qy-c+\partial_{\rm lim}g(y).
\]

This does not assert the invalid norm-gradient formula at zero. Given \textbf{any} complete resolvent output \(y\) at \(z\), choose \(w\in\partial_{\rm lim}g(y)\) with \(z-y=\lambda(Qy-c+w)\). Then

\[
(I+\lambda Q)(y-z)=\lambda(c-Qz-w),\qquad
\|y-z\|_2\le\lambda(\|Qz-c\|_2+L_p), \tag{D.9}\label{eq:dyn9}
\]

because \(Q\succ0\). The elementary bounds

\[
7/5<R<3/2,\quad \sqrt2<10/7,\quad 2^{1/14}<15/14
\]

give, for \(z\in\mathcal U\),

\[
\|Qz-c\|_2+L_p
\le\sqrt2(1-R/2)+2R/64+L_p
<\frac37+\frac3{64}+\frac52<3. \tag{D.10}\label{eq:dyn10}
\]

Therefore \textbf{every} complete output satisfies

\[
\|y-\bar x\|_2<1/64+3/128=5/128,
\quad |h_y|,|t_y|<5/(128\sqrt2)<1/32. \tag{D.11}\label{eq:dyn11}
\]

No localization of outputs was assumed in deriving \eqref{eq:dyn9}. In particular remote nonsmooth outputs and zero are excluded before using any smooth formula.

On the convex box \(\mathcal B\), both individual coordinates are at least \(15/16\). The usual norm Hessian formula gives

\[
0\preceq\nabla^2\|y\|_{7/4}
\preceq\tfrac34\|y\|_{7/4}^{-3/4}
\operatorname{diag}(y_i^{-1/4}).
\]

Here \(\|y\|_{7/4}\ge(15/16)R>1\) and
\((3/4)(16/15)^{1/4}<1\). Thus
\(\nabla^2\Phi(y)=Q-\nabla^2\|y\|_{7/4}\succ-I\).
The map \(y\mapsto y+\lambda\nabla\Phi(y)\) is strictly monotone on this convex box, because its Jacobian is at least \((1-\lambda)I\succ0\). All complete outputs from \eqref{eq:dyn11} lie in this box, so two such outputs at the same input must coincide.

For existence, \(\Phi\) is continuous and nonnegative; hence its proximal objective is continuous and coercive for every input and every positive step. A global minimizer exists, and Fermat places every global minimizer in the complete \(J_{\lambda F}\). The proved at-most-one output therefore gives exactly one complete output and exactly one global minimizer, and the two maps coincide.

Finally, write \(z=y+\lambda\nabla\Phi(y)\). By \eqref{eq:dyn8},

\[
\|z-\bar x\|_2^2
=\|y-\bar x\|_2^2+2\lambda D(h_y,t_y)
+\lambda^2\|\nabla\Phi(y)\|_2^2
\ge\|y-\bar x\|_2^2. \tag{D.12}\label{eq:dyn12}
\]

Thus \(y\in\mathcal U\), establishing the stated invariant ball for the complete resolvent and global proximal map.

\subsection{Actual PPA tracking and the exact slow-tail constant}

Fix \(0<\lambda\le1/128\). By complete coverage, uniqueness, and \eqref{eq:dyn12}, the trajectory starting at any \(x_0\in\mathcal U\) is defined uniquely for every integer \(k\ge0\) and stays in that ball. Summing \eqref{eq:dyn12} and using \eqref{eq:dyn8} gives

\[
\sum_{k=0}^{\infty}(h_{k+1}^2+t_{k+1}^4)<\infty.
\]

Therefore \(h_k,t_k\to0\), and \(x_k\to\bar x\). This argument uses the ordinary proximal equations, not a reduced scalar minimization algorithm.

The exact coordinate form of those equations is

\[
h_k=h_{k+1}+\frac\lambda2\varphi_h(h_{k+1},t_{k+1}),\qquad
t_k=t_{k+1}+\frac\lambda2\varphi_t(h_{k+1},t_{k+1}). \tag{D.13}\label{eq:dyn13}
\]

The factor \(1/2\) comes from the Euclidean metric
\(\|x-x_k\|_2^2=2[(h-h_k)^2+(t-t_k)^2]\).
From \eqref{eq:dyn1}--\eqref{eq:dyn3}, on the box,

\[
\begin{aligned}
\varphi_h/R&=4h+\tfrac38t^2+O(|h|t^2+t^4),\\
\varphi_t/R&=\tfrac34ht+\tfrac{11}{64}t^3
 +O(h^2|t|+|h||t|^3+|t|^5). \end{aligned}
\tag{D.14}\label{eq:dyn14}
\]

All remainders here are analytic with uniform bounds on the box; these are local asymptotic estimates derived from the explicit certificate, not assumptions about the trajectory.

If \(t\ne0\), its exact tangent derivative is

\[
\frac{\varphi_t}{Rt}=\frac34\frac h{s}
+\frac{11}{64}\frac{t^2}{s^3}-\frac{E'(t/s)}t.
\]

By \eqref{eq:dyn3}, \(|h|,|t|\le1/32\), and \(s\ge31/32\), the absolute value is at most

\[
\frac3{124}+\frac{11}{64\cdot32^2}(32/31)^3
+\frac{28}{32^4}(32/31)^5<\frac1{32}. \tag{D.15}\label{eq:dyn15}
\]

The bound extends continuously to \(t=0\). Since \(\lambda R/64<1\), the factor in the second equation of \eqref{eq:dyn13},
\(1+(\lambda/2)(\varphi_t/t)\), is positive everywhere in the box. Consequently \(t_{k+1}=0\) if and only if \(t_k=0\), and nonzero tangent coordinates keep their sign. On the axis \(t=0\), \eqref{eq:dyn1} gives \(\varphi=2Rh^2\), so \eqref{eq:dyn13} gives exactly the geometric formula in the theorem.

Now suppose \(t_0\ne0\). The tangent multiplier tends to one as \(h_k,t_k\to0\), and therefore

\[
\frac{t_k}{t_{k+1}}\longrightarrow1. \tag{D.16}\label{eq:dyn16}
\]

Set \(a=1+2\lambda R>1\), \(b_0=3\lambda R/16\), and \(q=1/a<1\). Introduce the analytic even function
\[
\kappa(u)=\frac{1-f(u)+u f'(u)}{u^2},\qquad \kappa(0)=3/8.
\]
The value at zero is its removable analytic extension. Exact differentiation
of \eqref{eq:dyn1} gives \(\varphi_h/R=4h+u^2\kappa(u)\).
For the actual output with \(s_{k+1}=1+h_{k+1}\), the first equation of
\eqref{eq:dyn13} therefore yields the exact tracking recurrence
\[
\begin{aligned}
w_{k+1}&=q\left(\frac{t_k}{t_{k+1}}\right)^2w_k
-q\frac{\lambda R}{2}\frac{\kappa(t_{k+1}/s_{k+1})}{s_{k+1}^2},\\
w_k&=h_k/t_k^2.
\end{aligned}
\tag{D.17}\label{eq:dyn17}
\]
No boundedness of \(w_k\) is presumed. By \eqref{eq:dyn16} and convergence
to the origin, the multiplier tends to \(q<1\) and the forcing tends to
\(-qb_0\). Choose any \(\theta\in(q,1)\); eventually the multiplier has
absolute value below \(\theta\) and the forcing is bounded. Iterating this
scalar inequality first proves \(w_k\) bounded. Subtract the fixed point
\(-qb_0/(1-q)\): the new forcing tends to zero, and the same contraction
then proves convergence. Hence
\[
\frac{h_k}{t_k^2}\to-\frac{qb_0}{1-q}=-\frac3{32}.
\tag{D.18}\label{eq:dyn18}
\]

This is the required radial-following proof. It is valid for all off-axis initial points in the certified ball; it neither assumes nor requires invariance of the radial critical curve.

Substitute \eqref{eq:dyn18} into the second equation of \eqref{eq:dyn13}, using \eqref{eq:dyn14} at the actual output. The remainders become \(o(t_{k+1}^2)\) after division by \(t_{k+1}\), and

\[
t_k=t_{k+1}\left[1+
\frac{13\lambda R}{256}t_{k+1}^2+o(t_{k+1}^2)\right]. \tag{D.19}\label{eq:dyn19}
\]

Writing \(\tau_k=|t_k|>0\), expansion of the reciprocal square gives

\[
\tau_{k+1}^{-2}-\tau_k^{-2}\longrightarrow
\frac{13\lambda R}{128}.
\]

Telescoping and taking Ces\`aro means yields
\(\tau_k^{-2}/k\to13\lambda R/128\), proving the tangent limit. Finally \eqref{eq:dyn18} implies \(h_k/t_k\to0\), and
\(\|x_k-\bar x\|_2=\sqrt2\sqrt{h_k^2+t_k^2}\sim\sqrt2\tau_k\).
This proves the stated exact distance limit.


\section{Interpretation and verification boundaries}
The structural transition concerns active curvature at every local minimum.
It does not imply that all stationary points are minima, that arbitrary
PPA initializations reach a minimum, or that a sparse ground-truth signal
is recovered. Nondegenerate local rates have their ordinary
strong-monotonicity interpretation. The quartic example also admits the
classical analytic gradient/KL interpretation. Here its complete-map
certification and exact actual-trajectory asymptotic distinguish the
algorithmic assertion from a reduced scalar curve calculation.
No exponent or tail is asserted uniformly across all upper-side models.
The step and basin constants are sufficient rather than optimal.

The accompanying Python checks verify exact finite rational identities,
analytic-majorant arithmetic, high-precision boundary expansions, and
actual two-variable implicit proximal equations. These corroborate the
analytic proofs, which carry the universal claims. The central
publication-priority gate remains subject to the two prioritized primary-text
gaps and the relevant secondary-body gap in Section 1.
An exact antecedent, if found, must be credited.
Renaming standard local consequences as PPA cannot replace lost novelty.

\begin{thebibliography}{9}
\small
\setlength{\itemsep}{3pt}
\bibitem{WangPenalty}
Y. Wang, \emph{New improved penalty methods for sparse reconstruction
based on difference of two norms}, author report, March 2013; revised
August 2014; Optimization Online upload, 2015.
\url{https://optimization-online.org/wp-content/uploads/2015/03/4849.pdf}.
\bibitem{Yin}
P. Yin, Y. Lou, Q. He, and J. Xin,
\emph{Minimization of $\ell_{1-2}$ for compressed sensing},
SIAM J. Sci. Comput. \textbf{37} (2015), A536--A563.
\href{https://doi.org/10.1137/140952363}{doi:10.1137/140952363}.
\bibitem{WangZhang}
D. Wang and Z. Zhang,
\emph{Generalized sparse recovery model and its neural dynamical
optimization method for compressed sensing},
Circuits Syst. Signal Process. \textbf{36} (2017), 4326--4353.
\href{https://doi.org/10.1007/s00034-017-0532-7}{doi:10.1007/s00034-017-0532-7}.
\bibitem{LouYan}
Y. Lou and M. Yan, \emph{Fast $L_1-L_2$ minimization via a proximal
operator}, J. Sci. Comput. \textbf{74} (2018), 767--785.
\href{https://arxiv.org/abs/1609.09530v4}{Author manuscript, arXiv:1609.09530v4}.
\bibitem{Huo}
L. Huo, W. Chen, H. Ge, and M. Ng,
\emph{$L_1-\beta L_q$ minimization for signal and image recovery},
SIAM J. Imaging Sci. \textbf{16} (2023), 1886--1928.
\href{https://doi.org/10.1137/22M1525363}{doi:10.1137/22M1525363}.
\bibitem{ZhaoNeural}
Y. Zhao, X. He, T. Huang, and J. Huang,
\emph{Smoothing inertial projection neural network for minimization
$L_{p-q}$ in sparse signal reconstruction},
Neural Networks \textbf{99} (2018), 31--41; published online
20 December 2017.
\href{https://doi.org/10.1016/j.neunet.2017.12.008}{doi:10.1016/j.neunet.2017.12.008}.
\bibitem{Fractional}
Y. Zhao, H. He, C. Ma, and H. Wang,
\emph{A unified fractional regularization framework for sparse recovery},
arXiv:2604.23184v2, 27 May 2026.
\url{https://arxiv.org/abs/2604.23184v2}.
\end{thebibliography}
\end{document}
